
Large language models can repair their own code when given execution feedback: run the generated solution, capture the output or exception, and feed it back to the model for correction. This repair loop paradigm, exemplified by Self-Debug \citep{Chen2023Teaching} and Reflexion \citep{Shinn2023Reflexion}, has become a standard technique for improving pass@k on benchmarks like HumanEval and MBPP. The dominant feedback signal is execution output --- a binary pass/fail plus stderr. The intuition is that more specific feedback should produce better repairs: if the model knows not just that it failed but why (a type error at line 5 with expected type \texttt{int}, received \texttt{str}), it should correct the error more reliably. Static analysis tools like mypy generate exactly this kind of structured, type-level diagnostic. Yet no controlled study has measured whether adding mypy to a repair loop produces measurable improvement over execution-only feedback.

The gap matters for a practical reason: execution feedback is cheap but coarse. mypy feedback is almost equally cheap (a subprocess call, approximately 30ms) but structured --- it provides error category, file location, and error message. If this structure translates into better repairs, the improvement is available to any practitioner today, requiring only a one-line addition to an existing repair loop. But without controlled evidence, we cannot answer: how much does it help, when does it help, and why does it help?

We conducted a five-sub-hypothesis controlled experiment to answer these questions. The experiment proceeds as a causal chain: we first verify that mypy produces actionable signal on our target benchmarks (h-e1), then verify the signal is consumed by the LLM (h-m1), then test whether the improvement is concentrated on type-error problems (h-m2), then measure the overall pass@1 improvement (h-m3), and finally test whether a formal verification extension (Z3 counterexample feedback) provides additional benefit (h-z1).

The results contain surprises at multiple levels. On HumanEval+ (164 problems, GPT-4o-mini, k=5 repair rounds, 3 seeds), execution+mypy repair achieves \textbf{87.20\% pass@1} versus \textbf{78.25\%} for execution-only repair --- a \textbf{+8.94 percentage point} improvement, consistent across all three seeds (seed 42: +6.7pp; seed 123: +11.0pp; seed 456: +9.2pp). mypy errors are eliminated within a single repair round (mean errors: 1.55 to 0.00 by round 2; Spearman $\rho = -0.707)$. So far, the story fits the intuition.

But the mechanism does not. We pre-registered a sub-hypothesis (h-m2) that the improvement would be concentrated on type-error problems --- the problems where mypy provides its diagnostic signal. It is not. Under k=5 repair rounds, type-error problems achieve \textbf{90.9\% repair rate under both conditions} (differential = 0.000). The improvement exists broadly across problem types. The most parsimonious interpretation --- though not yet confirmed via ablation (see L2, Section 6) --- is general diagnostic context enrichment: mypy output, regardless of whether it signals a type error, may augment the repair prompt with structured text that improves GPT-4o-mini's repair quality across all problem types. This interpretation requires a controlled ablation (Condition B-Null: equivalent prompt length but with non-mypy content) to confirm.

A second finding: mypy error rates differ dramatically between HumanEval+ (70\% of failing solutions have mypy errors) and MBPP+ (0\%). This benchmark divergence predicts where static analysis feedback will and will not help, and constrains all pass@1 claims to HumanEval+. The Z3 formal verification extension (h-z1) shows null improvement ($-0.9pp)$ on an arithmetic subset, constrained by low counterexample activation (16.1\%) and ceiling effects (base pass@1\_B = 86.6\%).

This paper makes the following contributions:

\begin{enumerate}
\item \textbf{First controlled empirical comparison} of execution+mypy repair versus execution-only repair on HumanEval+ (GPT-4o-mini, k=5, three seeds), showing a consistent +8.94pp pass@1 improvement.
\end{enumerate}

\begin{enumerate}
\item \textbf{Type-specificity refutation}: a pre-registered sub-hypothesis test establishing that the improvement is not concentrated on type-error problems (identical 90.9\% repair rates under both conditions at k=5), motivating the general context enrichment interpretation.
\end{enumerate}

\begin{enumerate}
\item \textbf{Benchmark-level mypy signal characterization}: HumanEval+ and MBPP+ have fundamentally different mypy error profiles (70\% vs 0\%), a structural difference that predicts where static analysis feedback will activate.
\end{enumerate}

\begin{enumerate}
\item \textbf{Z3 spec validation pipeline} achieving 91.1\% spec validity rate on arithmetic HumanEval+ problems, establishing feasibility for future formal verification work while identifying the ceiling and activation-rate limitations that currently limit its repair-loop utility.
\end{enumerate}

Section 2 reviews related work on LLM repair loops and static analysis feedback. Section 3 describes the experimental methodology and the five-sub-hypothesis causal chain. Section 4 details the experimental setup. Section 5 presents results across all sub-hypotheses. Section 6 discusses mechanistic implications, limitations, and boundary conditions. Section 7 concludes with future work directions.

